\documentclass{article}

% ── Track / review mode ────────────────────────────────────────────────────────
% Submission (double-blind, default):
\usepackage{neurips_2026}
% Camera-ready after acceptance — add "final":
%   \usepackage[main, final]{neurips_2026}
% Other tracks: [position], [eandd], [creativeai], [sglblindworkshop],
%   [dblblindworkshop], [education].  Add "final" for camera-ready in any track.
% Preprint / arXiv:
%   \usepackage[preprint]{neurips_2026}

% ── Packages ───────────────────────────────────────────────────────────────────
\usepackage[utf8]{inputenc}
\usepackage[T1]{fontenc}
\usepackage{hyperref}
\usepackage{url}
\usepackage{booktabs}
\usepackage{amsfonts}
\usepackage{amsmath,amssymb}
\usepackage{nicefrac}
\usepackage{microtype}
\usepackage{xcolor}
\usepackage{graphicx}

% ── Title ──────────────────────────────────────────────────────────────────────
\title{TITLE_GOES_HERE}

% Single author:
%   \author{Name\thanks{optional footnote} \\ Dept \\ University \\ \texttt{email}}
%
% Multiple authors — use \And (same line) or \AND (force new line):
\author{%
  First Author \\
  Department of Computer Science \\
  University Name \\
  City, Country \\
  \texttt{author@university.edu}
  \And
  Second Author \\
  Affiliation \\
  \texttt{email@affiliation.com}
}

% ── Document ───────────────────────────────────────────────────────────────────
\begin{document}

\maketitle

\begin{abstract}
% 1–2 paragraphs: problem statement, approach, key results, and significance.
% NeurIPS abstracts are typically 150–250 words.
\end{abstract}

\section{Introduction}
% Motivate the problem and state the research question.
% Summarise the contributions clearly (numbered list or bold bullets are common).
% End with a one-sentence roadmap of the paper.

\section{Related Work}
% Survey prior work grouped by theme or approach.
% Explain how this work differs from or improves on each group.

\section{Methodology}
% Describe the approach clearly enough for a reader to reproduce it.
% Include task formulation, model architecture, and algorithms.

\section{Experiments}
% Describe datasets, baselines, metrics, and implementation details.
% NeurIPS requires compute/resource information — state GPU type, training time, etc.

\section{Results}
% Report the main quantitative findings (tables and figures encouraged).
% Include ablations and discuss statistical significance where applicable.

\section{Discussion}
% Interpret the results: what do they imply beyond the numbers?
% Address unexpected findings and limitations of the experimental setup.

\section{Conclusion}
% Summarise the key contributions and findings in 1–2 paragraphs.
% State implications and suggest concrete directions for future work.

\section*{Limitations}
% Required by NeurIPS.  Use a starred (unnumbered) section.
% Cover scope, data coverage, compute requirements, evaluation constraints, etc.

% ── Bibliography ───────────────────────────────────────────────────────────────
\bibliographystyle{plainnat}
\bibliography{references}

% ── Appendix ───────────────────────────────────────────────────────────────────
\appendix

\section{Appendix}
\label{sec:appendix}
% Supplementary material: proofs, extended tables, additional experiments, etc.
% Appendix sections are numbered A, B, … automatically after \appendix.

% ── NeurIPS checklist (required, must be last) ─────────────────────────────────
\newpage
\input{checklist.tex}

\end{document}
